\section{Implementation Surface}
\label{sec:doctrine-implementation}

\subsection{Source Files}
\label{sec:doctrine-implementation-sources}

The doctrine corpus is implemented entirely as markdown source files within the
\texttt{doctrine/} directory of \texttt{\textasciitilde/process-intelligence}. This is a
deliberate architectural choice: the doctrine layer has no compiled source code, no build
artifacts, and no executable surface of its own. Its implementation is the specification
itself --- the 34 canonical files that constitute the authority from which downstream
implementations derive their type laws, algorithm boundaries, and authorization conditions.

The 25 files examined for this chapter cluster into five functional groups:

\paragraph{Upstream Law Files.} \texttt{BLUE\_RIVER\_DAM.md} (law stack, five maturity
levels, format covenant, the upstream dam principle),
\texttt{AUTONOMIC\_KNOWLEDGE\_ACTUATION.md} (MAPE-K loop specification with receipt
requirements for every actuation phase), \texttt{CONFORMANCE\_AS\_LAW.md} (token replay and
alignment-based conformance as process-level law enforcement, with type law assignments per
algorithm), \texttt{PROCESS\_INTELLIGENCE\_DEFINED.md} (canonical definition, five maturity
levels, upstream dam principle, anti-regression audit mesh specification).

\paragraph{Evidence and Receipt Law Files.} \texttt{EVIDENCE\_CHAIN.md} (typed one-way
lifecycle, \texttt{Evidence\textless{}T,State,W\textgreater{}} carrier, failure mode
taxonomy), \texttt{RECEIPT\_DOCTRINE.md} (Receipt\textless{}T,W\textgreater{} specification,
chain-of-custody law, board claim traceability, M\&A diligence surface),
\texttt{NAMED\_LAW\_REFUSAL.md} (named law inventory, anti-pattern prohibition).

\paragraph{Algorithm and Object Taxonomy Files.} \texttt{ALGORITHM\_TAXONOMY.md} (complete
fifteen-family taxonomy with compat/wasm4pm ownership assignments and gap status),
\texttt{FORMAL\_OBJECTS\_TAXONOMY.md} (type-theoretic specifications for all formal
objects), \texttt{OBJECT\_CENTRIC\_SUPREMACY.md} (OCEL 2.0 divergence/convergence
resolution), \texttt{VAN\_DER\_AALST\_CANON.md} (paper citation registry grounding all
algorithm families).

\paragraph{Lifecycle and Governance Files.} \texttt{PROCESS\_LIFECYCLE\_ONTOLOGY.md}
(twelve-stage lifecycle: Design $\to$ Simulation $\to$ Construction $\to$ Activation $\to$
Operation $\to$ Monitoring $\to$ Repair $\to$ Optimization $\to$ BoardProjection $\to$
Integration $\to$ Decommission $\to$ Archive), \texttt{GRADUATION\_LAW.md}
(\texttt{GraduationCandidate\textless{}T\textgreater{}} specification),
\texttt{MA\_READY\_PROCESS\_INTELLIGENCE.md} (six board-admissibility criteria, M\&A deck
manufacturing algorithm), \texttt{PROCESS\_TRUTH\_AUTHORITY.md} (authority stack and
refusal capital definition).

\paragraph{Authorization and Gap Files.} \texttt{DOWNSTREAM\_AUTHORIZATION\_LAW.md} (six
authorized downstream workflows receipted against ALIVE\_001),
\texttt{RESEARCH\_AUTHORITY.md} (authority table and GAP\_001 finding),
\texttt{full-lifecycle-process.md}, \texttt{spr\_thesis\_actuation.md},
\texttt{lattice-monotonicity-verification.md}, \texttt{lifecycle\_algorithms.md},
\texttt{bpmn-or-join-completion.md}, \texttt{wasm-boundary-law.md}.

\subsection{Commands}
\label{sec:doctrine-implementation-commands}

The doctrine corpus defines no runnable commands of its own. It is the authority layer, not
the execution layer. However, \texttt{DOWNSTREAM\_AUTHORIZATION\_LAW.md} specifies the
verification command that must be run before initiating any authorized downstream workflow:

\begin{verbatim}
cat ~/process-intelligence/checkpoints/PROCESS_INTELLIGENCE_ALIVE_001.md
\end{verbatim}

This command verifies the seal status of the ALIVE\_001 checkpoint and is the precondition
for all six authorized downstream workflows. It is the operational translation of the
authorization gate logic:

\[
  \text{DownstreamRefactorAuthorized}(r, w) \iff
  \text{ALIVE\_001.sealed} = \top \;\land\; w \in \text{AuthorizedWorkflows} \;\land\;
  r \in w.\text{targets}
\]

The anti-regression audit mesh described in \texttt{PROCESS\_INTELLIGENCE\_DEFINED.md}
references an \texttt{audit\_crown\_gate\_all.sh} script that coordinates 22 subordinate
audit scripts. This script is specified in doctrine but its implementation resides in the
downstream wasm4pm-compat codebase.

\subsection{Data and Ontology Surfaces}
\label{sec:doctrine-implementation-data}

The doctrine corpus defines several data-level ontology surfaces that downstream
implementations must conform to.

\paragraph{Process Lifecycle Ontology.} The twelve-stage lifecycle is an ordered ontology
with lawful transitions. Each stage produces a specific evidence artifact: Design produces
the formal process model; Monitoring produces fitness and precision scores; Repair produces
non-conformance remediation plans; Archive produces the final closure receipt. Stages may not
be skipped; the lifecycle is a total order over stage identifiers.

\paragraph{Algorithm Taxonomy Data.} The fifteen-family algorithm taxonomy in
\texttt{ALGORITHM\_TAXONOMY.md} serves as the data surface for the manufacturing surface
combinatorial formula:
\[
  \text{Surface} = |S| \times |L| \times |A| \times |F| \times |R| \times |Q| \times |B| \times |D|
\]
where $S$ is the algorithm family set, $L$ is the log type set, $A$ is the admissibility
state set, $F$ is the fitness metric family set, $R$ is the refusal reason set, $Q$ is the
quality metric set, $B$ is the board claim type set, and $D$ is the lifecycle stage set.

\paragraph{Paper Citation Registry.} Every algorithm family entry in the taxonomy includes
explicit paper citations with author, year, and algorithm name. This registry is the data
layer that the \texttt{ggen} manufacturing machinery uses to manufacture type-law surfaces
from published papers. Citations are not decorative; they are the grounding authority for
every type shape in the compat layer.

\subsection{Templates and Rendered Artifacts}
\label{sec:doctrine-implementation-templates}

The \texttt{PROCESS\_INTELLIGENCE\_SPR\_THESIS.md} file is the primary rendered artifact of
the doctrine corpus. It is a Sparse Priming Representation (SPR) of the full thesis argument,
manufactured at the ALIVE\_001 seal point and explicitly stating that it ``was generated at
ALIVE\_001 seal.'' It compresses the full argument into a machine-readable priming surface
that downstream language model invocations may use to reconstruct the thesis context without
re-reading all 34 doctrine files.

The \texttt{DOWNSTREAM\_AUTHORIZATION\_LAW.md} file functions as a rendered authorization
template: it specifies six workflow templates, each with a named target repository, a named
prompt file path, a scope description, a gate condition, and a receipt requirement. These
templates are the manufacturing surface for authorized downstream work.

The \texttt{spr\_thesis\_actuation.md} file documents the actuation of the SPR thesis ---
the process by which the SPR was manufactured and its receipt chain established. This file
is the receipt for the template itself, demonstrating that the doctrine corpus applies its own
evidence requirements recursively to its own manufacturing outputs.
